% Opening of the old chapter c28, its \chapter command removed; generated
\section*{Truncations of the beta lattice and Hermite coefficients}\label{chap:tr}

Section~\ref{ch:rz4} truncates the lattice of Biane, Pitman and Yor for zeta: the truncated zeros move at the first order rate $\kappa(\rho)/n$ (Proposition~\ref{rz:prop:firstorder}), the side of approach is tied to the Riemann Siegel phase and so to Gram's law (Theorem~\ref{rz:thm:kappaexact}), the zero $\rho_{213}$, approached from the wrong side, refutes the one sided conjecture (Corollary~\ref{rz:cor:Tfalse}), and symmetrising by the functional equation leaves a statement that is not refuted (Proposition~\ref{rz:prop:symtrunc}). We carry this programme to the Dirichlet beta function $L(s,\chi_4)$, with the lattice $\lambda_k=2\pi^2(k-\frac12)^2$ and the completed $L$ function as Mellin transform, and add the orders $m\ge1$ of cumulant matching, an exact mean corrected truncation as a finite sum of incomplete gamma functions, and the drift at all $379$ zeros below height $500$. The second half is a separate subject: the expansion of $\Xi$ in Hermite functions at a small scale, whose coefficients sample $\Xi$ at the turning points of the Hermite functions, up to an amplitude factor and an error $O(a^4)$, and have been used to count zeros from the P\'olya kernel alone.

The chapter is based on \cite{P26trunc,P26hermite}. Where a recomputation differed from a note the recomputed value is printed, and a statement not recomputed is marked unchecked. Three statements of the notes are restated: the drift law needs a hypothesis on the first unmatched cumulant of the remainder (Theorem~\ref{tr:thm:drift}, Proposition~\ref{tr:prop:gamma}), the exclusion of one sided criteria is conditional on one zero with a negative drift (Corollary~\ref{tr:cor:nogo}), and the asymptotic of the second note is proved in an additive form, uniform in $m$, the relative form failing at the zeros of $\Xi$ (Theorem~\ref{tr:thm:herm}, Proposition~\ref{tr:prop:zeros}).

Notation. $\xi_4(s)=\xi(s,\chi_4)$ is that of Section~\ref{sec:dbeta:beta}, and $\Lambda$ is the de Bruijn Newman constant. The variable $\omega$ is the exit time of Remark~\ref{rem:dbeta:exit} divided by $2\pi$. $\vartheta_\beta$ and $Z_\beta$ are the Riemann Siegel and Hardy functions of $\chi_4$, so that $\vartheta$ and $Z$ remain those of zeta. $\Phi$ is the kernel of \eqref{eq:ch2:Phi}, $H_t$ the de Bruijn Newman function of Section~\ref{ch:dbn}, and $H_n$ with an integer index the Hermite polynomial of Section~\ref{ht:sec:jensen}. The letter $a$ is the Hermite scale in the second half, and $\lambda_*=2\pi^2$ in the first.

Floating point computations are called computed and never certified; the zero counts of this chapter are of that kind. The results are these.
\begin{enumerate}
\item GRH for $\chi_4$ is equivalent to a local zero free condition on the truncated Mellin transforms (Theorem~\ref{tr:thm:local}, proved). The truncated zeros move by $\delta_mn^{-(2m+1)}$ for the hard cutoff and the mean correction, and by a multiple of the first unmatched cumulant difference in general (Theorem~\ref{tr:thm:drift}, Corollary~\ref{tr:cor:named}, proved); the formula for $m\ge2$ with a generalised gamma convolution remainder is false (Proposition~\ref{tr:prop:gamma}, proved), and the drift of the first zero is computed for three kinds of remainder (Table~\ref{tr:tab:drift}).
\item The sign of the real part of the drift is a Gram parity of $\chi_4$ with two corrections (Theorem~\ref{tr:thm:phase}, proved). Among the $379$ zeros below height $500$ the mean correction approaches from the left exactly at the five zeros with anomalous Gram parity (computed, Table~\ref{tr:tab:gram}). A zero with negative drift excludes a one sided criterion at that order (Corollary~\ref{tr:cor:nogo}, proved); the premise is computed, not proved.
\item The truncations have transient zeros in the left half strip (Table~\ref{tr:tab:counts}, computed). The symmetrised truncations satisfy a Hurwitz criterion (Theorem~\ref{tr:thm:symhurwitz}, proved) and keep every simple critical zero on the line locally (Proposition~\ref{tr:prop:symlocal}, proved); beyond that only the heights of the computed tables are reached.
\item The Hermite coefficients of $\Xi$ at scale $a$ are the Taylor coefficients of $e^{z^2/2}$ times the backward heat flow of $\Xi$ for the time $2a^2$ (Theorem~\ref{tr:thm:heat}, proved). They sample $\Xi$ at the turning points $a\sqrt{4m+1}$ with an additive error $O(a^4)$ uniform in $m$ (Theorem~\ref{tr:thm:herm}, proved), and the first negative coefficient reads off $\gamma_1$ (Corollary~\ref{tr:cor:first}, proved; Table~\ref{tr:tab:first}, computed).
\item The number of sign changes of the coefficients equals the number of zeros of zeta up to height $1000$ at scales up to about $5$ to $6$, with a lag of order one at the top edge (Remark~\ref{tr:obs:count}, computed); what a verified count would prove is Proposition~\ref{tr:prop:verified}.
\end{enumerate}
